\section*{ID9808: Relation: 153600/π ∼ 4.88×10\textasciicircum{}4}
\paragraph{Metadata.}
PiID: 2067235850200 | RTEID: RTE:P38086.6421:T6.2064 | UUID: 4667d8bd-8d82-5086-b00f-35a02a614d55

\paragraph{Canonical Statement.}
es as there are kissing sphere contacts. This sequence – effectively a kissing-number-weighted CPMG sequence – flattens out environmental noise up to an extremely high order. The electron spin coherence time \$T\_2\$ is enhanced roughly by a factor of \$N\{\textbackslash{}max\}/\textbackslash{}pi\$, and our decoupling analysis predicts that the noise Hamiltonian’s norm is suppressed below the system Hamiltonian by a factor \$1/N\_\{\textbackslash{}max\}\textasciicircum{}\{\textbackslash{}pi\}\$. In practical terms, an electron spin that normally loses phase coherence in microseconds could maintain coherence for seconds or longer by using a 153,600-length \$\textbackslash{}pi\$-pulse sequence (since \$153600/\textbackslash{}pi \textbackslash{}sim 4.88\textbackslash{}times10\textasciicircum{}4\$). We return to this point when discussing our metamaterial’s ability to keep qubits coherent at room temperature. We have similarly formulated protection protocols for Cooper pairs (superconducting qubits), magnons (quantum spin waves), phonons (mechanical quanta), and polaritons (light-matter hybrids). Each protocol introduces \$\textbackslash{}pi\$ or kissing-number dependencies in the relevant Hamiltonian terms to harness topological symmetry. For a Cooper pair in a transmon qubit, for instance, the Josephso

\paragraph{Canonical Equation.}
\[
153600/\pi \sim 4.88\times10^4
\]

\paragraph{Dictionary.}
Relation: 153600/π ∼ 4.88×10\textasciicircum{}4 — 153600/π ∼ 4.88×10\textasciicircum{}4

\paragraph{Encyclopedia.}
Relation: 153600/π ∼ 4.88×10\textasciicircum{}4 is a scaling / order-of-magnitude relation, written 153600/π ∼ 4.88×10\textasciicircum{}4. It expresses a scaling or proportionality, capturing how the quantity grows with its parameters. It is built from π. Within the Trawin Zero Point registry it serves as one element of the framework's mathematical narrative.
